% Results from the 16 validated matched 600 messages/s trials.
% This passage and the figures are prepared for review before online inclusion.
The following figures compare the metrics in Table~\ref{tab:metric-roles}
using equal-weight means of the 2 runs for each response and traffic
concentration. Bands and vertical whiskers show the range between the
runs, rather than confidence intervals. Time-series means require valid
values from both runs, so unavailable observations remain gaps. Gray
dotted lines mark the scheduled action at minute~1, while colored dotted
lines mark each run's confirmed activation. The dashed throughput line
marks the continuing input rate of 600~messages/s.

Figure~\ref{fig:skew-mean-control} shows sustained backlog growth and
increasing message age without intervention in both conditions. With
12 high-traffic partitions, redistribution reduced backlog and confirmed
recovery in both runs. With 4 high-traffic partitions, its initial
reduction was followed by continued accumulation, and the ages of
waiting messages increased. Lag skew remained elevated even after some
responses reduced backlog to a low level, so its magnitude alone does
not establish whether further intervention is needed.

\begin{figure}[htbp]
\centering
\includegraphics[width=\linewidth,height=.78\textheight,keepaspectratio]{figures/skew-mean-control.pdf}
\caption{Runtime control signals for the two traffic concentrations.}
\label{fig:skew-mean-control}
\end{figure}

The measured arrival rates in Figure~\ref{fig:skew-mean-candidates}
show approximately 40~messages/s per high-traffic partition in the
12-partition condition and 120~messages/s in the 4-partition condition.
The latter therefore places more demand on each partition even though
the total input remains unchanged. An audit of the saved records found eligible busy intervals for
Consumer~0 in both conditions and for Consumers~4 and~5 in the
4-partition condition. Their effective rates were approximately
259--260, 267, and 293~messages/s, respectively. These capacity points
pool completed attempts over observed time in eligible intervals,
while their whiskers show the observed interval-rate range. They are
retrospective diagnostics from the same trials rather than independent
calibration or means of 2 run-level capacity estimates. Consumers~1,
2, and~3 had no qualifying interval, and the balanced reference did
not establish a capacity estimate for any consumer.

\begin{figure}[htbp]
\centering
\includegraphics[width=\linewidth]{figures/skew-mean-candidates.pdf}
\caption{Pre-action arrival rates and retrospective busy-period capacity estimates.}
\label{fig:skew-mean-candidates}
\end{figure}

Figure~\ref{fig:skew-mean-performance} shows how completion rates and
latency changed during evaluation. Throughput includes distinct
warm-up messages completed during evaluation, which explains why
recovery peaks exceed the continuing input rate. Window p99 uses
evaluation messages completed in each 30-second interval, retaining
their original producer timestamps. The run-level outcomes in
Figure~\ref{fig:skew-mean-outcomes} instead follow the evaluation cohort
through the drain cutoff, and the p99 points average the 2 run-level
percentiles without pooling messages.

\begin{figure}[htbp]
\centering
\includegraphics[width=\linewidth]{figures/skew-mean-performance.pdf}
\caption{Throughput and window p99 during evaluation.}
\label{fig:skew-mean-performance}
\end{figure}

With 12 high-traffic partitions, redistribution completed every
evaluation message and recovered in a mean of 341~seconds after the
scheduled action. Targeted scaling recovered sooner, at 200~seconds,
but increased mean run-level p99 from 86.6 to 113.3~seconds. With
4 high-traffic partitions, redistribution retained a mean of 1,570
evaluation messages at the cutoff, or 0.44\%, whereas targeted scaling
completed the cohort and recovered in both runs. Its mean run-level
p99 was also lower, at 112.1 rather than 124.4~seconds. The benefit of
adding consumers therefore differed between the two traffic
concentrations under the tested partition maps. The measured
partition rates imply about 279--282~messages/s for Consumer~0 after
4-partition redistribution, above its observed busy-period range of
255--266~messages/s. With 12 high-traffic partitions, its corresponding
input was about 201--205~messages/s. This comparison is consistent
with the different backlog trajectories, but it does not establish
prospective prediction accuracy because the service estimates use
these same trials.

Normal Kafka scaling achieved the lowest mean run-level p99 in both
conditions, at 65.0 and 77.1~seconds, respectively. However, neither
4-partition run met the recovery criterion during continuing input,
despite completing every evaluation message by the drain cutoff.
Both final Kafka assignments placed 2 high-traffic partitions on at
least 1 consumer, whereas targeted scaling distributed them across
4 consumers. In this condition, Kafka scaling also had a higher mean
SLO violation rate than targeted scaling, at 54.2\% rather than
31.7\%. These outcomes show why consumer count and p99 alone do not
fully describe the effectiveness of a response.

\begin{figure}[htbp]
\centering
\includegraphics[width=\linewidth,height=.78\textheight,keepaspectratio]{figures/skew-mean-outcomes.pdf}
\caption{Evaluation-message outcomes and confirmed recovery times.}
\label{fig:skew-mean-outcomes}
\end{figure}

Figure~\ref{fig:skew-mean-costs} shows that redistribution used less
memory than either scaling response in both conditions. For the
4-partition workload, redistribution and targeted scaling had similar
CPU costs, about 19.4 and 19.6~core-minutes, although only targeted
scaling completed all evaluation messages and confirmed recovery.
Targeted scaling took longer to activate than the other interventions,
so its eventual recovery benefit must be considered alongside that
delay. Resource costs cover observed consumer CPU and RSS memory over
evaluation and drain, with more than 99\% measurement coverage.
Matched placement reduces one source of variation, but 2 repetitions
on shared nodes do not establish performance guarantees beyond these
conditions.

\begin{figure}[htbp]
\centering
\includegraphics[width=\linewidth,height=.70\textheight,keepaspectratio]{figures/skew-mean-costs.pdf}
\caption{Consumer resource cost and decision-to-active time.}
\label{fig:skew-mean-costs}
\end{figure}
\FloatBarrier
